\section{Diseño Experimental Propuesto}

Dado que el Marco de Aprendizaje Inclusivo y Adaptativo (MAIA) constituye una
arquitectura de referencia y no un sistema desplegado, su validación se plantea
como un protocolo escalonado que progresa desde la evaluación aislada de cada
módulo hasta la validación del sistema integrado en el aula. Esta gradación es
coherente con el diseño desacoplado por fusión tardía descrito en la Propuesta,
y permite distinguir con claridad qué propiedades pueden verificarse de forma
temprana sobre datos etiquetados y cuáles requieren un despliegue real.

% ------------------------------------------------------------
\subsection{Hipótesis de Validación}
% ------------------------------------------------------------

El desempeño esperado de MAIA se expresa mediante tres hipótesis falsables, una
por cada módulo central de la arquitectura:

\begin{itemize}
    \item \textbf{H1 (Seguimiento cognitivo).} El clasificador cognitivo,
    basado en un modelo BERT ajustado finamente, alcanza una precisión igual o
    superior al 85\,\% en los seis niveles de la Taxonomía de Bloom sobre un
    corpus curricular costarricense en español, tomando como referencia el
    83\,\% y 89\,\% reportados por Das et al.~\cite{Das2020} en inglés.

    \item \textbf{H2 (Percepción afectiva).} El módulo de percepción afectiva
    multimodal identifica los estados de frustración, aburrimiento y atención
    con una precisión igual o superior al 80\,\% y una latencia máxima de
    200\,ms por fotograma sobre el hardware objetivo, umbral dentro del cual la
    retroalimentación se percibe como continua~\cite{AISTY2024}.

    \item \textbf{H3 (Control docente).} El sistema opera de forma autónoma en
    al menos el 70\,\% de las sesiones, entendiendo por autonomía aquellas que
    concluyen sin superar el umbral de impacto configurable que activa la
    intervención del docente.
\end{itemize}

% ------------------------------------------------------------
\subsection{Fase 1: Validación de Componentes}
% ------------------------------------------------------------

La primera fase evalúa cada módulo de manera aislada y sin participación de
estudiantes, sobre conjuntos de datos etiquetados, lo que permite iterar sobre
los modelos sin riesgo ético. El módulo de seguimiento cognitivo se evalúa sobre
un corpus de ítems del currículo del Ministerio de Educación Pública (MEP)
anotados por nivel de Bloom; para garantizar la calidad de las etiquetas se
propone una doble anotación por especialistas curriculares y la medición del
acuerdo entre anotadores mediante el coeficiente kappa de
Cohen~\cite{Cohen1960}. El desempeño se cuantifica mediante la precisión global
y la medida F1 promediada por clase (F1 macro), apropiada ante el desbalance
esperado entre los seis niveles de Bloom, complementadas con la matriz de
confusión para identificar confusiones sistemáticas entre niveles cognitivos
adyacentes. Como línea base de comparación se emplean clasificadores clásicos de
menor complejidad, replicando el enfoque de Das et al.~\cite{Das2020}. El módulo
de percepción afectiva se valida sobre un repositorio público de expresiones
faciales en contexto educativo, tomando como línea base una configuración basada
únicamente en visión que se contrasta con la configuración multimodal completa,
de modo que se justifique empíricamente la estrategia de fusión tardía. Para
este módulo se reportan la precisión y la medida F1 de cada estado afectivo,
junto con la latencia de inferencia por fotograma ---expresada en percentiles,
mediana y percentil 95--- sobre el hardware objetivo, con el fin de verificar el
umbral de 200\,ms establecido en la hipótesis H2. En ambos módulos las
particiones de entrenamiento y prueba se separan por estudiante para evitar
fugas de identidad.

% ------------------------------------------------------------
\subsection{Fase 2: Validación de Integración}
% ------------------------------------------------------------

La segunda fase examina el sistema completo en sesiones controladas, con el
objetivo de medir la latencia de extremo a extremo, la proporción de sesiones
autónomas (hipótesis H3), el número de intervenciones docentes por sesión y la
recompensa acumulada del agente de refuerzo como indicador de la estabilidad del
bucle de adaptación con control docente. La línea base de esta fase es la
comparación entre el motor de adaptación operando como aprendizaje por refuerzo
puro frente al motor inicializado con la política de reglas pedagógicas del MEP,
con el fin de evidenciar que dicha política mitiga el problema de arranque en
frío descrito en la Propuesta.

% ------------------------------------------------------------
\subsection{Fase 3: Estudio Ecológico en Aulas Integradas}
% ------------------------------------------------------------

La fase final consiste en un cuasi-experimento en aulas reales que contrasta el
uso de MAIA frente a una condición de instrucción no adaptativa. Las variables
dependientes incluyen los indicadores de logro curricular, el compromiso del
estudiante, la carga de trabajo percibida por el docente ---medida con el Índice
de Carga de Tarea de la NASA (NASA Task Load Index,
NASA-TLX)~\cite{HartStaveland1988}--- y la usabilidad del sistema ---evaluada con
la Escala de Usabilidad del Sistema (System Usability Scale,
SUS)~\cite{Brooke1996}. Dado que el tamaño de muestra esperado en este contexto
es reducido, el análisis estadístico privilegia pruebas no paramétricas y el
reporte del tamaño del efecto por encima de la mera significancia estadística.

% ------------------------------------------------------------
\subsection{Consideraciones Éticas y Amenazas a la Validez}
% ------------------------------------------------------------

Por involucrar la participación de seres humanos en la investigación, y de
manera especialmente sensible por tratarse de una población vulnerable, el
protocolo requiere aprobación de un comité de ética, consentimiento informado de
los responsables y asentimiento de los estudiantes, y el procesamiento de datos sensibles en el borde
(\textit{edge computing}) conforme a la Ley N\textsuperscript{o}\,8968. Entre las
principales amenazas a la validez se anticipan el efecto de novedad, el sesgo del
docente al aprobar o rechazar decisiones ---que retroalimenta directamente al
motor de adaptación---, el posible sesgo del modelo afectivo entre estudiantes
con distintos perfiles de expresividad, y la limitada generalización derivada de
la heterogeneidad de hardware entre instituciones.
